\documentclass[a4paper,12pt]{article}

\usepackage[spanish]{babel}
\usepackage[utf8]{inputenc}
\usepackage{enumerate}
\usepackage{amsmath}
\usepackage{extsizes}
\usepackage{amssymb}
\usepackage{dsfont}
\usepackage{graphicx}
\usepackage{cancel}
\usepackage[usenames]{color}
\usepackage[dvipsnames]{xcolor}
\usepackage{accents}
\usepackage{flushend}
\usepackage{tikz}
\usepackage[LGR,T1]{fontenc}
\newcommand{\textgreek}[1]{\begingroup\fontencoding{LGR}\selectfont#1\endgroup}
\usetikzlibrary{arrows,automata}
\usepackage{multicol}
\setlength{\columnsep}{1cm}
\usepackage{listings}
\usepackage{graphics,graphicx, float} %para incluir imágenes y colocarlas



\usepackage[hidelinks]{hyperref}

\usepackage[vmargin=3cm,hmargin=3cm]{geometry}
%\setlength\parindent{0pt}

\setlength\parindent{0pt}


% Carpeta con las imágenes
%\graphicspath{{}}

\begin{document}


	\begin{center}
		\LARGE{\textbf{Visión por Computador (2018-2019)} \\ Grado en Ingeniería Informática y Matemáticas \\ Universidad de Granada }
		\vspace*{2.5cm}

		\rule{\textwidth}{1.6pt}\vspace*{-\baselineskip}\vspace*{4pt}
		\rule{\textwidth}{1.6pt}\vspace*{-\baselineskip}\vspace*{2pt}
		\vspace{0.5cm}

		\Huge{Descriptores HOG en la detección de peatones}

		\vspace{0.5cm}
		\rule{\textwidth}{1.6pt}\vspace*{-\baselineskip}\vspace*{2pt}
		\rule{\textwidth}{1.6pt}\vspace*{-\baselineskip}\vspace*{4pt}

		\vspace{4cm}

\begin{figure}[h!]
	\centering
	\includegraphics[scale=0.75]{./Imagenes/etsiit.jpeg}
	\label{fig:logoETSIIT}
\end{figure}

		\vspace{4cm}
		\large{Ignacio Aguilera Martos y Diego Asterio de Zaballa \\ \today }

	\end{center}




\newpage

\tableofcontents

\newpage

%%%%%%%%%%%%%%%%%%%%%%%%%%%%%%%%%%%%%%%%%%%%%%%%%%%%%%%%%%%%%%%%%%%%%%%%%%%%%%%%%%%%%%
%%             Descripción del problema y el enfoque que le hemos dado              %%
%%%%%%%%%%%%%%%%%%%%%%%%%%%%%%%%%%%%%%%%%%%%%%%%%%%%%%%%%%%%%%%%%%%%%%%%%%%%%%%%%%%%%%

\section{Descripción del problema y enfoque de la resolución}

\subsection{Problema a resolver}

\subsection{Descriptores HOG}

El problema mencionado anteriormente se ha enfrentado desde
distintos puntos de vista lo que ha dado lugar a diversos
descriptores. Nuestra implementacion se basa en el paper de
Navneet Dalal y Bill Triggs en el que desarrollan un descriptor
basado en el gradiente de una imagen llamado descriptor HOG.
El metodo se basa en calcular los gradientes de una imagen y a
continuacion construir histogramas locales basado en la
orientacion de dichos gradientes.

La idea basica sobre la que se desarrolla la tecnica es que las
caracteristicas de forma y apariencia de un objeto de forma local
se pueden resumir gracias a la distribucion de la orientacion de
los gradientes de la imagen.

A continuacion se describe la implementacion de la tecnica a
grandes rasgos. En primer lugar la ventana de la imagen sobre la
que se va a construir el vector de caracteristicas se divide en
celdas. Estas celdas son pequenas regiones de la imagen. En cada
una de estas celdas se calcula el histograma de la direccion de
los gradientes que en ella se encuentran. Esto define la re
A continuacion estas
celdas se agrupan en bloques para que sea posible normalizarlas.


diEl descriptor HOG es un vector de caracteristicas 
La construccion del descriptor HOG se apoya fundamentalmente en el gradiente de la imagen.

\subsection{Fases de la resolución}

%%%%%%%%%%%%%%%%%%%%%%%%%%%%%%%%%%%%%%%%%%%%%%%%%%%%%%%%%%%%%%%%%%%%%%%%%%%%%%%%%%%%%%
%%                           Valoración de los resultados                           %%
%%%%%%%%%%%%%%%%%%%%%%%%%%%%%%%%%%%%%%%%%%%%%%%%%%%%%%%%%%%%%%%%%%%%%%%%%%%%%%%%%%%%%%

\section{Valoración de los resultados}

%%%%%%%%%%%%%%%%%%%%%%%%%%%%%%%%%%%%%%%%%%%%%%%%%%%%%%%%%%%%%%%%%%%%%%%%%%%%%%%%%%%%%%
%%                               Propuestas de mejora                               %%
%%%%%%%%%%%%%%%%%%%%%%%%%%%%%%%%%%%%%%%%%%%%%%%%%%%%%%%%%%%%%%%%%%%%%%%%%%%%%%%%%%%%%%

\section{Trabajo futuro: propuestas de mejora}

\normalsize


% Bibliografía.
%-----------------------------------------------------------------
%\onecolumn
%\bibliography{referencias}
%\bibliographystyle{plain}
%\nocite{*}

\end{document}
